\section{Protocole entre la montre et le téléphone}
\label{sec:protocole}

\subsection{Principes}

\begin{itemize}
  \item La montre \emph{demande}, le téléphone \emph{répond}~; hors liaison et état d'attente, le
        téléphone n'envoie rien de lui-même.
  \item Un message ne contient qu'un élément (une pièce, une rangée d'historique, une alerte), jamais
        une liste. Une liste est une tête suivie de ses éléments (C6).
  \item Toute demande porte un numéro de séquence que la réponse reprend. La montre ignore les
        réponses d'une séquence qui n'est plus la sienne (l'utilisateur a changé d'écran).
  \item Les températures voyagent en dixièmes de degré, dans l'unité du compte (DEC-12). Les
        durées en minutes, les heures en minutes depuis minuit.
  \item Les textes ne voyagent que lorsqu'ils varient (noms de pièces, abréviations)~; tout le reste
        est un code que la montre traduit (DEC-12).
  \item La taille des tampons se règle sur le plus gros message (la pièce, environ 70 octets avec
        deux textes)~; la valeur est à confirmer sur chaque plateforme (point ouvert ci-dessous).
\end{itemize}

\pointouvert{Quelle taille donner aux tampons AppMessage sur chaque plateforme (le SDK fournit la taille maximale par plateforme)~? Le plus gros message (une pièce, environ 70 octets) doit y tenir avec marge.}

\subsection{Messages}

\begin{table}[H]
  \centering\small
  \begin{tabularx}{\linewidth}{@{}l>{\raggedright\arraybackslash}p{23mm}X@{}}
    \toprule
    \textbf{Réf.} & \textbf{Nom} & \textbf{Contenu} \\
    \midrule
    \multicolumn{3}{@{}l}{\textit{De la montre vers le téléphone}} \\
    MSG-01 & HELLO & Langue de la montre, forme de l'écran, nombre de colonnes d'historique souhaité. Envoyé à l'ouverture, après \texttt{ready}. \\
    MSG-02 & GET & Séquence~; objet demandé (pièces, détail, historique, confort, alertes, présence, raccourcis, plage de consigne)~; indice de pièce~; indicateur \enquote{actualiser}. \\
    MSG-03 & DO & Séquence~; commande (forcer, programmation, éteindre, présence, raccourci, lier, délier)~; indice de pièce~; température en dixièmes~; durée (code)~; numéro de raccourci. \\
    \midrule
    \multicolumn{3}{@{}l}{\textit{Du téléphone vers la montre}} \\
    MSG-11 & STATUS & Séquence~; état de la liaison (non liée, perdue, active)~; âge de la donnée affichée (secondes)~; indicateurs~: hors connexion, attente~; mode et état de présence~; nombre d'alertes. \\
    MSG-12 & ROOMS\_HEAD & Séquence~; nombre de pièces~; unité (C ou F). \\
    MSG-13 & ROOM & Séquence~; indice~; nom (20 octets)~; abréviation (4)~; température mesurée~; consigne~; puissance (\%)~; tranche de couleur~; indicateurs (chauffe, forcée, éteinte, hors ligne, fenêtre ouverte)~; code de mode et durée restante. \\
    MSG-14 & DETAIL & Séquence~; indice~; code de mode~; durée restante d'un forçage minuté~; heure de fin (absente pour \enquote{jusqu'à nouvel ordre})~; heure et température du prochain changement~; humidité (\%)~; minutes de fenêtre ouverte. \\
    MSG-15 & HISTORY\_ROW & Séquence~; indice~; niveaux des $n$ colonnes (2 bits chacune)~; masque des colonnes inconnues~; indice de la colonne en cours. \\
    MSG-16 & COMFORT & Séquence~; indice~; température et humidité en dixièmes~; code du verdict~; indicateur \enquote{humidité non mesurée}. \\
    MSG-17 & ALERT & Séquence~; indice~; indice de pièce~; nature~; détail (minutes ou code). \\
    MSG-18 & PRESENCE & Séquence~; mode (Présent, Absent, Auto)~; état~; nombre d'appareils suivis par la géolocalisation et nombre de ceux qui sont à la maison~; indicateur \enquote{Auto indisponible} (mémorisé par le téléphone après un refus). \\
    MSG-19 & SHORTCUTS & Séquence~; masque des raccourcis offerts. \\
    MSG-20 & CAPS & Séquence~; indice de pièce~; minimum, maximum et pas de consigne, en dixièmes. \\
    MSG-21 & ACK & Séquence~; résultat (pris en compte, refusé, échec, indisponible, hors connexion, Auto indisponible). Suivi de la vue relue. \\
    MSG-22 & LINK\_CODE & Code à saisir (6 caractères)~; secondes restantes~; état (en attente, validé, expiré, refusé). \\
    \bottomrule
  \end{tabularx}
  \caption{Messages}
\end{table}

\subsection{Échanges types}

\paragraph{Ouverture.} MSG-01~; la watchapp affiche sa copie (DEC-07)~; MSG-02 \enquote{pièces}~;
MSG-11~; MSG-12~; MSG-13 $\times n$~; la watchapp remplace sa copie lorsque les $n$ pièces sont
reçues, ou signale un résultat partiel après un délai (§\ref{sec:erreurs}).

\paragraph{Commande.} MSG-03~; MSG-21 (résultat)~; la vue relue (MSG-13 ou MSG-12 puis MSG-13 $\times n$)~;
la montre n'affiche l'état de la pièce qu'à cette réception (DEC-10).

\paragraph{Historique.} MSG-02 \enquote{historique}~; MSG-11~; MSG-15 $\times n$ (une par pièce). La montre
dessine chaque rangée à sa réception, ce qui montre une grille qui se remplit plutôt qu'un
écran vide.

\subsection{Vues}

Le tableau~\ref{tab:vues} donne la correspondance entre les écrans et les vues que
le pkjs construit (J-11).

\begin{table}[H]
  \centering\small
  \begin{tabularx}{\linewidth}{@{}>{\raggedright\arraybackslash}p{30mm}>{\raggedright\arraybackslash}p{32mm}X@{}}
    \toprule
    \textbf{Écran} & \textbf{Messages} & \textbf{Données d'origine (Tado)} \\
    \midrule
    1 Pièces & MSG-11, 12, 13 & \texttt{zoneStates}, \texttt{state}. \\
    2 Détail & MSG-14 & \texttt{zoneStates} (ou relecture d'une pièce). \\
    3 Pilotage & MSG-14 & Mode et consigne en cours~; dernier forçage mémorisé par le pkjs. \\
    4 Consigne & MSG-20 & \texttt{capabilities}. \\
    5 Durée & aucun & Constantes de la montre. \\
    6 Menu & MSG-11 & Nombre d'alertes, présence. \\
    7 Historique & MSG-15 & \texttt{dayReport} par pièce. \\
    8 Présence & MSG-18 & \texttt{state}, \texttt{mobileDevices}. \\
    9 Raccourcis & MSG-19 & \texttt{zoneStates}, \texttt{state}. \\
    10 Bien-être & MSG-16 & \texttt{airComfort}, \texttt{zoneStates}. \\
    11 Alertes & MSG-17 & \texttt{zones}, \texttt{zoneStates}. \\
    12 Liaison & MSG-22 & \texttt{login.tado.com}. \\
    \bottomrule
  \end{tabularx}
  \caption{Vues et données d'origine}
  \label{tab:vues}
\end{table}
